% =====================================================================
%  Hojun Kang — Curriculum Vitae
%
%  내용은 이 파일이 아니라  ../data/profile.yml  과 Notion DB에 있습니다.
%      python ../scripts/build.py    →  _generated.tex 갱신
%      latexmk cv.tex                →  cv.pdf   (latexmkrc 가 XeLaTeX 지정)
%  이 파일은 서식(폰트·여백·레이아웃)만 담당합니다.
% =====================================================================
\documentclass[11pt,a4paper]{article}

% ---- 폰트: Times New Roman, 한글은 바탕 --------------------------------
\usepackage{fontspec}
\setmainfont{Times New Roman}

\usepackage[a4paper,top=0.85in,bottom=0.85in,left=0.9in,right=0.9in,
            footskip=0.45in]{geometry}
\usepackage{titlesec}
\usepackage{needspace}
\usepackage{fancyhdr}
\usepackage{lastpage}
\usepackage[hidelinks]{hyperref}

\hypersetup{colorlinks=false, pdfborder={0 0 0},
            pdftitle={Hojun Kang -- Curriculum Vitae}, pdfauthor={Hojun Kang}}

\setlength{\parindent}{0pt}
\setlength{\parskip}{0pt}
\linespread{1.0}
\frenchspacing
\raggedbottom

% ---- 꼬리말 ------------------------------------------------------------
\pagestyle{fancy}
\fancyhf{}
\renewcommand{\headrulewidth}{0pt}
\fancyfoot[L]{\footnotesize Hojun Kang}
\fancyfoot[R]{\footnotesize \thepage\ of \pageref{LastPage}}

% ---- 섹션: 대문자 볼드 + 가는 선 ---------------------------------------
\titleformat{\section}
  {\normalsize\bfseries\MakeUppercase}
  {}{0pt}{}
  [{\vspace{0.15em}\hrule height 0.4pt}]
\titlespacing*{\section}{0pt}{1.15em plus 0.25em minus 0.2em}{0.55em}

\titleformat{\subsection}[block]
  {\needspace{4.5\baselineskip}\normalsize\itshape}{}{0pt}{}
\titlespacing*{\subsection}{0pt}{0.7em}{0.22em}

% =====================================================================
%  엔트리 매크로
% =====================================================================

% \entry{기간}{제목}{기관}{비고}
%  좌우 두 칸 모두 같은 크기로 두고 \strut 으로 첫 줄 baseline 을 맞춘다.
\newcommand{\entry}[4]{%
  \needspace{2.4\baselineskip}%
  \noindent
  \begin{minipage}[t]{0.235\textwidth}%
    \raggedright\strut #1%
  \end{minipage}%
  \hfill
  \begin{minipage}[t]{0.745\textwidth}%
    \raggedright\strut\textbf{#2}%
    \ifx\relax#3\relax\else\\ #3\fi
    \ifx\relax#4\relax\else\\[0.1em]{\small #4}\fi
  \end{minipage}%
  \par\vspace{0.5em}%
}

% \pub{...}  — APA 행잉 인덴트 참고문헌 한 건
\newcommand{\pub}[1]{%
  \needspace{2\baselineskip}%
  \par\hangindent=1.4em\hangafter=1
  #1\par\vspace{0.34em}%
}

% \me{...}     — 본인 이름 강조
\newcommand{\me}[1]{\textbf{#1}}
% \status{...} — 심사 상태. 색·배경 없이 이탤릭 괄호로만 적는다.
\newcommand{\status}[1]{\textit{(#1)}}
\newcommand{\statusok}[1]{\textit{(#1)}}

% =====================================================================
\begin{document}

% ---- 헤더 -------------------------------------------------------------
\begin{center}
  {\LARGE\bfseries Hojun Kang}\\[0.2em]
  {\small Division of Business Administration, Dongduk Women's University}\\[0.15em]
  {\small Humanities Building Room 406, 60 Hwarang-ro 13-gil, Seongbuk-gu, Seoul 02748, Korea}\\[0.15em]
  {\small \href{mailto:hjkang@dongduk.ac.kr}{hjkang@dongduk.ac.kr} \quad $\cdot$ \quad +82 10-9770-6275}
\end{center}

\vspace{0.9em}
\hrule height 0.4pt
\vspace{0.4em}

% ---- 자동 생성 본문 ----------------------------------------------------
% 이 파일은 scripts/build.py 가 생성합니다. 직접 고치지 마세요.

\section{Education}
\entry{Mar 2020 -- Aug 2023}{Ph.D., Management Information Systems}{Sogang University, Seoul, Korea}{Advisor: Sang-Gun Lee\newline Dissertation: Cryptocurrency, Brave New Realm or Mere Dream? Using the Dunning-Kruger Effect and Size Effect for Empirical Analysis}

\section{Academic Appointments}
\entry{Sep 2026 -- Present}{Assistant Professor, Division of Business Administration}{Dongduk Women's University}{Teaching and research in management information systems, with emphasis on AI applications to business and financial-market analysis\newline Research program spanning AI governance, graph-based financial risk detection, and the economics of digital assets}
\entry{Apr 2025 -- Aug 2026}{Senior Researcher, AI Policy Research Division}{Software Policy \& Research Institute (SPRI)}{Ethical AI implementation and policy framework development\newline Lead analyst of AI regulatory trends, industry impact assessment, and sustainable AI ecosystem policy recommendations}
\entry{Sep 2024 -- Apr 2025}{Research Professor, BK21 Business Education Research Group}{Sogang University}{Research report writing, teaching, and project monitoring for the BK21 program}
\entry{May 2024 -- Aug 2024}{Researcher, BK21 Business Education Research Group}{Sogang University}{BK21 project proposal development and progress reporting}
\entry{Sep 2023 -- Aug 2024}{Honorary Professor, Business School}{Sogang University}{Development and submission of the BK21 project proposal; overall management and coordination}
\entry{Mar 2018 -- Dec 2022}{Research Assistant, ITRC Intelligent Blockchain Research Center}{Sogang University}{Economic research and analysis of blockchain technology}

\section{Publications}
\subsection{Refereed journal articles}
\pub{\me{Kang, H.}, \& Lee, S.-G. (2026). Frozen Markets, Falling Prices: Illiquidity and Bounded Price Discovery in U.S. Neighborhood Housing Markets. \emph{The Journal of Real Estate Finance and Economics}. \status{SSCI}}
\pub{\me{Kang, H.}, Hwang, K., Lee, Y.-K., \& Lee, S.-G. (2026). Did Spot Crypto ETFs Reallocate Venue Liquidity? Fragile Evidence. \emph{Finance Research Letters}, 111, 110744. \status{SSCI Q1}}
\pub{\me{Kang, H.}, \& Lee, S.-G. (2026). TRACE: Transaction Representation Learning for Large-Scale Ethereum Phishing Detection under Limited Labels. \emph{Decision Support Systems}. \status{SCIE Q1}}
\pub{\me{Kang, H.}, \& Lee, S.-G. (2026). Quieter Markets, Unchanged Risks: Frequency Connectedness and the Spot Bitcoin ETF. \emph{International Review of Economics and Finance}. \status{SSCI Q1}}
\pub{\me{Kang, H.}, Kwon, M., \& Lee, S.-G. (2026). Declining Connectedness, Rising Influence: TVP-VAR Evidence from the Magnificent Seven. \emph{Finance Research Letters}. \status{SSCI Q1}}
\pub{\me{Kang, H.}, \& Hwang, K. (2025). FORTRESS: Fraud-oriented Transformer with Random Traversal for Ethereum Security Surveillance. \emph{Information Sciences}. \status{SCIE Q1}}
\pub{Lee, S., Lee, S.-G., \& \me{Kang, H.} (2025). Unleashing STO Potential: An AHP-Based Analysis of Critical Success Factors. \emph{Blockchain: Research and Applications}. \status{SCIE Q1}}
\pub{\me{Kang, H.}, \& Lee, S.-G. (2025). Market Phases and Price Discovery in NFTs: A Deep Learning Approach to Digital Asset Valuation. \emph{Journal of Theoretical and Applied Electronic Commerce Research}. \status{SSCI Q1}}
\pub{Hwang, K., \& \me{Kang, H.} (2025). The Role of Household Savings in Fiscal and Monetary Policy: A DSGE Analysis of China. \emph{Asian and African Studies}. \status{SSCI Q2}}
\pub{Liu, H., Lee, S.-G., \& \me{Kang, H.} (2025). Graph-based Phishing Detection on Ethereum: A Comparative Analysis of GNN and Traditional Embedding Methods. \emph{Journal of Information Systems}. \status{KCI}}
\pub{\me{Kang, H.}, Lee, S.-G., \& Park, S.-Y. (2022). Information Efficiency in the Cryptocurrency Market: The Efficient-Market Hypothesis. \emph{Journal of Computer Information Systems}. \status{SCIE Q2}}
\pub{\me{Kang, H.}, Lee, S.-G., Park, S.-Y., \& Lee, H.-H. (2019). A Comparative Study on ICT Industries over South Korea, the US, and Japan Using Industry Relations Table. \emph{The Korean Academy of Business History}, 34(2), 89--107. \status{KCI}}
\pub{\me{Kang, H.}, \& Park, B. (2018). An Analysis on Customers' Perception of Korean Food in the US Using Text Mining. \emph{Logos Management Review}, 16(4), 155--168. \status{KCI}}
\subsection{Books}
\pub{Lee, S.-G., Kim, Y. J., Kim, A., \& \me{Kang, H.} (2021). Data Analytics Methodology That Anyone Can Use. Sigma Press.}

\section{Working Papers}
\subsection{Revise and resubmit}
\pub{\me{Kang, H.}, Lee, H.-J., \& Lee, S.-G. (2026). The Asymmetric Effect of Spot Crypto ETP Approvals on Equity-Market Integration: Synthetic Difference-in-Differences Evidence. \emph{International Review of Financial Analysis}. \status{Major revision}}
\pub{\me{Kang, H.}, \& Lee, S.-G. (2026). Guilt by Association: Interpretable Multiplex Graph Learning of Accounting-Risk Contagion. \emph{Information Sciences}. \status{Resubmitted after minor revision}}
\pub{Byun, H., Lee, S.-G., \& \me{Kang, H.} (2026). The Paradox of Regulation: Why Does Housing Supply Fail to Respond to Price Signals? \emph{Journal of Real Estate Research}. \status{Resubmitted after major revision}}
\pub{\me{Kang, H.}, \& Lee, S.-G. (2025). Structural-Relational Dual Encoding for Blockchain Fraud Detection: The Walk2Attn Framework. \emph{Journal of Supercomputing}. \status{Major revision}}
\pub{\me{Kang, H.}, Lee, S.-G., \& Trimi, S. (2026). Unmasking the Hidden Threat: A Hierarchical Multi-Scale Graph Convolutional Network for Detecting Ethereum Phishing. \emph{Information Sciences}. \status{Minor revision}}
\subsection{Under review}
\pub{\me{Kang, H.}, \& Lee, S.-G. (2026). Freezing Digital Dollars: What the Stop Button Captures and What It Costs. \emph{Journal of Finance}.}
\pub{\me{Kang, H.}, \& Lee, S.-G. (2026). Exit Without Discipline: Delegation, Attention, and Accountability in Blockchain Governance. \emph{Management Science}.}
\pub{\me{Kang, H.}, \& Lee, S.-G. (2026). The Market for Governance Votes. \emph{The Review of Financial Studies}.}
\pub{\me{Kang, H.}, \& Lee, S.-G. (2026). Mechanism, Agent, or Harness? A Validity Boundary for LLM-Agent Economic Simulation. \emph{Information Systems Research}.}
\subsection{Work in progress}
\pub{Park, W., Lee, S.-G., \& \me{Kang, H.} Who Bears the Repricing? Segmented Incidence of a Capital-Driven Cycle in Korean Logistics Real Estate.}
\pub{\me{Kang, H.}, \& Lee, S.-G. Selective realization in tokenized real estate: portfolio-level evidence.}
\pub{\me{Kang, H.}, \& Lee, S.-G. The Production of Market Efficiency.}
\pub{\me{Kang, H.}, \& Lee, S.-G. Stranded Dollars: The Private Production of Transferability.}
\pub{\me{Kang, H.}, \& Lee, S.-G. (2026). Priced Out at the Threshold: Mandatory Insurance Costs, Debt-to-Income Limits, and Access to Mortgage Credit.}

\section{Conference Presentations}
\pub{\me{Kang, H.}, \& Lee, S.-G. (Jun 2026). MATE: Multi-Actor Graph Neural Networks for Corporate Delisting Prediction. KMIS, Gyeongju, Korea.}
\pub{\me{Kang, H.}, Lee, S.-G., \& Lee, D. W. (Jun 2025). The Collision of Technology and Reputation: Information Asymmetry and Investor Bias in Cryptocurrency Markets. KMIS, Seoul, Korea.}
\pub{\me{Kang, H.}, \& Lee, S.-G. (Oct 2024). Unmasking the Hidden Threat: Detecting Ethereum Phishing Nodes with Graph-Based Node2Vec-Transformer for Scam Detection Model. DSI, Arizona, USA.}
\pub{\me{Kang, H.}, \& Lee, S.-G. (Nov 2024). Gas Price Stabilization in Ethereum 2.0. KMIS, Seoul, Korea.}
\pub{\me{Kang, H.}, \& Lee, S.-G. (Oct 2024). Detecting Ethereum Phishing Nodes in NT4SD. E-Bridge, Seoul, Korea.}
\pub{\me{Kang, H.}, \& Lee, S.-G. (Aug 2024). Scanning Ethereum Phishing Node. KASBA, Gyeongju, Korea.}
\pub{\me{Kang, H.}, \& Lee, S.-G. (May 2024). Information Asymmetry in Cryptocurrency Markets: Evidence from the Dunning-Kruger Effect. KMIS, Seoul, Korea.}
\pub{\me{Kang, H.}, \& Lee, S.-G. (Jun 2023). The Dunning-Kruger Effect and Information Disparities in Cryptocurrency Trading Behavior. KMIS, Jeju, Korea.}
\pub{\me{Kang, H.}, Park, S.-Y., \& Lee, S.-G. (Jul 2021). Covid-19 Simulation Study with Beacon System: Through Bayesian Methods and Epidemic Models. KMIS, Yeosu, Korea.}
\pub{\me{Kang, H.}, Park, S.-Y., \& Lee, S.-G. (Dec 2020). Information Efficiency in the Cryptocurrency Market: The Efficient Market Hypothesis. KMIS, Seoul, Korea.}
\pub{\me{Kang, H.}, Park, S., Kim, J., \& Park, B. (Dec 2017). Analysis of Success Factors of FPS Games: Using Text Mining Analysis. Korean Intelligent Information Systems Society, Seoul, Korea.}
\pub{\me{Kang, H.}, Ko, Y., Park, H., Yoo, J., \& Park, B. (Nov 2017). An Analysis on Customers' Perception of Korean Food in the US Using Text Mining. Korea Society of Electronic Commerce, Mokpo, Korea.}

\section{Funded Projects}
\entry{Apr 2026 -- Present}{High-Impact AI Impact Assessment Project}{Software Policy \& Research Institute (SPRI)}{}
\entry{Apr 2025 -- Present}{MSIT-based AI Trend Monitoring}{Software Policy \& Research Institute (SPRI)}{}
\entry{Apr 2025 -- Present}{Survey on the Status of the AI Industry}{Software Policy \& Research Institute (SPRI)}{}
\entry{Apr 2025 -- Mar 2026}{Research on Improving the AI Industry Classification System}{Software Policy \& Research Institute (SPRI)}{}
\entry{Apr 2025 -- Mar 2026}{Research on AI Adoption in SMEs}{Software Policy \& Research Institute (SPRI)}{}
\entry{Apr 2025 -- Mar 2026}{Strategies for Promoting Data Interoperability and the Right to Data Portability in the Age of AI}{Software Policy \& Research Institute (SPRI)}{}
\entry{Oct 2023 -- Apr 2025}{ITRC Web3.0 Technology Research Center}{Sogang University}{}
\entry{Oct 2024 -- Feb 2025}{Forecasting NFT Market Dynamics using Transformer-based Deep Learning: A Predictive Analysis}{Korea Association of Game Industry (KGames)}{}
\entry{Sep 2024 -- Dec 2024}{Global Digital Talent Attraction: Policy Development and Implementation Strategies}{Software Policy \& Research Institute (SPRI)}{}
\entry{Jan 2024 -- Jul 2024}{Development and Implementation of Distributed Storage Systems for Energy Big Data Analytics}{Ministry of Trade, Industry and Energy (MOTIE)}{}
\entry{Feb 2023 -- Jul 2023}{Chungnam Digital \& Game Industry Development Feasibility Study and Operation Strategy}{Chungnam Information \& Culture Industry Promotion Agency (CTIA)}{}
\entry{Sep 2022 -- Dec 2022}{Virtual Power Plant (VPP) Economic Feasibility Study}{Korea Energy Agency (KEA)}{}
\entry{Sep 2021 -- Jan 2022}{Central Bank Digital Currency (CBDC) and Tax Administration Research}{Korea Institute of Public Finance (KIPF)}{}
\entry{Oct 2020 -- Dec 2020}{Digital New Deal --- Busan City}{Office of National Assembly Member Byeong-su Seo}{}

\section{Honors and Awards}
\entry{Oct 2024}{Best Paper Award}{E-Bridge Fall Conference}{}
\entry{Jun 2021}{Outstanding Case Study Award}{Korea Society of Management Information Systems (KMIS) Spring Conference}{}
\entry{Dec 2020}{Best Paper Award}{KMIS Fall Conference}{}
\entry{Nov 2017}{Best Paper Award}{Korea Society of Electronic Commerce Fall Conference}{}


\end{document}
